\documentclass[a4paper,12pt]{article}

\usepackage[ngerman]{babel}
\usepackage[utf8]{inputenc}
\usepackage{geometry}
\usepackage{listings}
\usepackage{xcolor}

\geometry{margin=2.5cm}

\title{Praktische Übung: CQRS und Event Sourcing mit Node.js}
\author{Modul 321}
\date{\today}

\lstset{
    basicstyle=\ttfamily\small,
    breaklines=true,
    frame=single,
    backgroundcolor=\color{gray!10}
}

\begin{document}

\maketitle

\section*{Ziel der Übung}

In dieser Übung entwickelt ihr ein einfaches Bankkonto-System und wendet dabei zwei wichtige Architekturkonzepte an:

\begin{itemize}
\item \textbf{CQRS (Command Query Responsibility Segregation)}
\item \textbf{Event Sourcing}
\end{itemize}

\textbf{Wichtig:} Der Kontostand darf niemals direkt gespeichert werden!

Stattdessen werden alle Änderungen als \textbf{Events} gespeichert.

\section*{Funktionalität}

Euer System soll:

\begin{itemize}
\item Konto eröffnen
\item Geld einzahlen
\item Geld abheben
\item Kontostand anzeigen
\item Alle Events anzeigen
\item Zustand durch Events rekonstruieren
\end{itemize}

\section*{Projektstruktur}

\begin{lstlisting}
cqrs-event-sourcing-bank/
|
|-- data/
|   `-- events.jsonl
|
`-- src/
    |-- eventStore.js
    |-- account.js
    |-- commands.js
    |-- queries.js
    `-- app.js
\end{lstlisting}

\section*{JSONL Format}

Jede Zeile ist ein Event:

\begin{lstlisting}
{"type":"AccountOpened","accountId":"A-100","owner":"Mia"}
{"type":"MoneyDeposited","accountId":"A-100","amount":100}
\end{lstlisting}

\section*{Aufgabe 1: Event Store}

Implementiert:

\begin{itemize}
\item appendEvent(event)
\item readEvents()
\end{itemize}

\section*{Aufgabe 2: Zustand berechnen}

Startzustand:

\begin{lstlisting}
{
  accountId: null,
  owner: null,
  balance: 0,
  isOpen: false
}
\end{lstlisting}

Implementiert:

\begin{itemize}
\item applyEvent(account, event)
\item rebuildAccountFromEvents(events, accountId)
\end{itemize}

\section*{Aufgabe 3: Commands}

Implementiert:

\begin{itemize}
\item openAccount(accountId, owner)
\item depositMoney(accountId, amount)
\item withdrawMoney(accountId, amount)
\item closeAccount(accountId)
\end{itemize}

\textbf{Wichtig:}  
Commands verändern niemals direkt den Zustand!

Sie erzeugen nur Events.

\section*{Aufgabe 4: Queries}

Implementiert:

\begin{itemize}
\item getAccount(accountId)
\item getAllEvents()
\end{itemize}

\textbf{Wichtig:}  
Queries verändern keine Daten.

\section*{Aufgabe 5: Testprogramm}

Beispiel:

\begin{lstlisting}
openAccount("A-100", "Mia");
depositMoney("A-100", 200);
withdrawMoney("A-100", 50);

console.log(getAccount("A-100"));
\end{lstlisting}

\section*{Beobachtungsaufgaben}

Beantwortet:

\begin{enumerate}
\item Wird der Kontostand direkt gespeichert?
\item Welche Events entstehen?
\item Kann man den Zustand aus Events rekonstruieren?
\item Was passiert bei mehrfacher Ausführung?
\end{enumerate}

\section*{CQRS verstehen}

\textbf{Command-Seite:}
\begin{itemize}
\item openAccount
\item depositMoney
\item withdrawMoney
\item closeAccount
\end{itemize}

\textbf{Query-Seite:}
\begin{itemize}
\item getAccount
\item getAllEvents
\end{itemize}

\section*{Erweiterung: Read Model}

Erstellt eine Datei:

\begin{lstlisting}
accounts_read_model.json
\end{lstlisting}

Diese darf den Kontostand speichern, wird aber aus Events berechnet.

\section*{Reflexionsfragen}

\begin{enumerate}
\item Warum wird der Kontostand nicht direkt gespeichert?
\item Was ist der Vorteil von Event Sourcing?
\item Was ist ein Nachteil?
\item Warum trennt man Commands und Queries?
\item Was passiert, wenn das Read Model gelöscht wird?
\item Kann man es wiederherstellen?
\end{enumerate}

\section*{Zusatzaufgaben}

\subsection*{Mehrere Konten}

Erstellt mehrere Konten und zeigt sie an.

\subsection*{Überweisung}

Implementiert:

\begin{lstlisting}
transferMoney(from, to, amount)
\end{lstlisting}

\subsection*{REST API}

Optional mit Express:

\begin{lstlisting}
POST /accounts
POST /accounts/:id/deposit
GET /accounts/:id
\end{lstlisting}

\section*{Bewertung}

\begin{itemize}
\item Events werden korrekt gespeichert
\item Trennung von Commands und Queries
\item Zustand wird aus Events berechnet
\item Fehler werden behandelt
\item Struktur ist verständlich
\end{itemize}

\section*{Diskussion}

\begin{itemize}
\item Wann ist Event Sourcing sinnvoll?
\item Wann ist es zu kompliziert?
\item Welche Systeme könnten davon profitieren?
\end{itemize}

\end{document}
